% ============================================================
%  CONFIGURATION DES POLICES
% ============================================================
% Objectif: rendu fiable en francais (accents, ligatures, ponctuation)

\usepackage{unicode-math}
\defaultfontfeatures{Ligatures=TeX,Scale=MatchLowercase}

% --- Police principale (serif) ---
% Rendu "rapport classique" proche du style de reference.
\IfFontExistsTF{TeX Gyre Pagella}{
  \setmainfont[Numbers=Proportional]{TeX Gyre Pagella}
}{
  \setmainfont[Numbers=Proportional]{Libertinus Serif}
}

% --- Sans-serif ---
\IfFontExistsTF{Noto Sans}{
  \setsansfont{Noto Sans}
}{
  \IfFontExistsTF{Source Sans 3}{
    \setsansfont{Source Sans 3}
  }{
    \setsansfont{TeX Gyre Heros}
  }
}

% --- Monospace ---
\IfFontExistsTF{JetBrains Mono}{
  \setmonofont[Scale=0.92]{JetBrains Mono}
}{
  \IfFontExistsTF{Noto Sans Mono}{
    \setmonofont[Scale=0.94]{Noto Sans Mono}
  }{
    \setmonofont[Scale=0.95]{TeX Gyre Cursor}
  }
}

% --- Math ---
\IfFontExistsTF{TeX Gyre Pagella Math}{
  \setmathfont{TeX Gyre Pagella Math}
}{
  \setmathfont{Libertinus Math}
}

% ========== Familles optionnelles (si présentes) ==========
% Serif alternatives (Windows)
\IfFontExistsTF{Georgia}{
  \newfontfamily\GeorgiaFamily{Georgia}[Ligatures=TeX]
}{
  \newfontfamily\GeorgiaFamily{TeX Gyre Pagella}[Ligatures=TeX]
}

\IfFontExistsTF{Cambria}{
  \newfontfamily\CambriaFamily{Cambria}[Ligatures=TeX]
}{
  \newfontfamily\CambriaFamily{TeX Gyre Pagella}[Ligatures=TeX]
}

% Sans-serif alternatives (Windows)
\IfFontExistsTF{Arial}{
  \newfontfamily\ArialFamily{Arial}[Ligatures=TeX]
}{
  \newfontfamily\ArialFamily{TeX Gyre Heros}[Ligatures=TeX]
}

\IfFontExistsTF{Corbel}{
  \newfontfamily\CorbellFamily{Corbel}[Ligatures=TeX]
}{
  \newfontfamily\CorbellFamily{TeX Gyre Heros}[Ligatures=TeX]
}

% Monospace alternative (Windows)
\IfFontExistsTF{Consolas}{
  \newfontfamily\ConsolasFam{Consolas}[Scale=0.9]
}{
  \newfontfamily\ConsolasFam{TeX Gyre Cursor}[Scale=0.95]
}

% ========== Commandes pratiques ==========
\newcommand{\HeadingFamily}{\sffamily}
\newcommand{\BodyFamily}{\rmfamily}
\newcommand{\CodeFamily}{\ttfamily}

\newcommand{\arial}[1]{{\ArialFamily #1}}
\newcommand{\corbel}[1]{{\CorbellFamily #1}}
\newcommand{\georgia}[1]{{\GeorgiaFamily #1}}
\newcommand{\cambria}[1]{{\CambriaFamily #1}}
\newcommand{\consolastext}[1]{{\ConsolasFam #1}}

\newcommand{\boldtext}[1]{\textbf{#1}}
\newcommand{\boldserif}[1]{{\CambriaFamily\textbf{#1}}}
\newcommand{\smallcapstext}[1]{\textsc{#1}}

% Attention: \emp dépend de HeigBlue (défini dans colors.tex)
\newcommand{\emp}[1]{\textit{\textcolor{HeigBlue}{#1}}}

\newcommand{\citationtext}[1]{{\GeorgiaFamily\itshape\small«~#1~»}}
